\section*{Hoofdstuk \ref{ch:b3:cell-cycle-apoptosis} --- Sturing van de celcyclus en geprogrammeerde celdood}

\begin{solution}{exo:b3:cell-cycle-apoptosis:1}
G1, S, G2, M. Voortgang door G1 en het \omterm{def:b3:cell-cycle-apoptosis:checkpoints}{restrictiepunt}: \omterm{def:b3:cell-cycle-apoptosis:cdk}{cycline}
D--Cdk4/6; G1/S: \omterm{def:b3:cell-cycle-apoptosis:cdk}{cycline} E--Cdk2; S: \omterm{def:b3:cell-cycle-apoptosis:cdk}{cycline} A--Cdk2; G2/M: \omterm{def:b3:cell-cycle-apoptosis:cdk}{cycline}
B--Cdk1. Het anafase-bevorderende complex (APC/C) vernietigt \omterm{def:b3:cell-cycle-apoptosis:cdk}{cycline} B
en securine en beëindigt de mitose.
\end{solution}

\begin{solution}{exo:b3:cell-cycle-apoptosis:2}
Hartwell: de \emph{cdc}-mutanten van bakkersgist, in vastgelegde stadia
stilgezet, en het begrip Start. Nurse: \emph{cdc2} van splijtgist als
het kinase dat de mitose timet, en zijn menselijke homoloog (Cdk1) die
de gist redt --- conservering. Hunt: \omterm{def:b3:cell-cycle-apoptosis:cdk}{cycline} in zee-egeleieren, een
eiwit dat de hele cyclus door wordt gemaakt en bij elke deling wordt
vernietigd.
\end{solution}

\begin{solution}{exo:b3:cell-cycle-apoptosis:3}
\omterm{def:b3:cell-cycle-apoptosis:checkpoints}{Restrictiepunt} (laat in G1): groeifactoren, voedingsstoffen, grootte;
grijpt in op \omterm{def:b3:cell-cycle-apoptosis:cdk}{cycline} D--Cdk4/6 en Rb. G2/M: DNA intact en
gerepliceerd; \omterm{def:b3:dna-repair:ddr}{ATM} en ATR remmen Cdc25, zodat Cdk1 gefosforyleerd en
inactief blijft. De spoelfiguur: elke kinetochoor gehecht; ongehechte
kinetochoren maken Mad2--Cdc20-complexen die het APC/C remmen.
\end{solution}

\begin{solution}{exo:b3:cell-cycle-apoptosis:4}
\omterm{def:b3:cell-cycle-apoptosis:apoptosis}{Apoptose}: de cel krimpt, het \omterm{def:b3:chromatin-epigenetics:nucleosome}{chromatine} verdicht zich, het DNA wordt in
een ladder geknipt, het membraan blaast op maar blijft gesloten,
\omterm{def:b3:cell-cycle-apoptosis:apoptosis}{fosfatidylserine} komt bloot, het lijk wordt verzwolgen, geen ontsteking.
\omterm{def:b3:cell-cycle-apoptosis:apoptosis}{Necrose}: de cel zwelt op, het membraan scheurt, de inhoud lekt weg, het
DNA wordt willekeurig afgebroken, ontsteking. De beulen: \omterm{def:b3:cell-cycle-apoptosis:apoptosis}{caspasen},
cysteïneproteasen die achter aspartaat knippen.
\end{solution}

\begin{solution}{exo:b3:cell-cycle-apoptosis:5}
Stijging van $5$ naar $30$ met $1$ per minuut: \qty{25}{min}. Afbraak van
$30$ naar $5$ met halveringstijd \qty{2}{min}: $\log_{2}(30/5)\times 2 =
\qty{5.2}{min}$. Periode ongeveer \qty{30}{min}, mitose \qty{17}{\%} daarvan.
\end{solution}

\begin{solution}{exo:b3:cell-cycle-apoptosis:6}
(a) $C$ kan nooit onder $C_{1}$ zakken: vastgezet in de mitose op de
hoge tak. (b) Geen remmende fosforylering: de drempel $C_{2}$ ligt
lager, de mitose begint vroeg bij een geringe grootte (het
``wee''-fenotype). (c) Cdk1 blijft gefosforyleerd: de hoge tak is
onbereikbaar, stilstand in G2 met verlengde cellen. (d) Als bij (b), en
het \omterm{def:b3:cell-cycle-apoptosis:checkpoints}{controlepunt} G2/M, dat via Wee1 en Cdc25 werkt, kan de cel niet
langer tegenhouden.
\end{solution}

\begin{solution}{exo:b3:cell-cycle-apoptosis:7}
G1-kern in cytoplasma van de S-fase: haar oorsprongen zijn vergund en
het cytoplasma levert werkzaam \omterm{def:b3:cell-cycle-apoptosis:cdk}{cycline} E/A--Cdk2, dus gaat zij meteen de
S-fase in. G2-kern in cytoplasma van de S-fase: haar oorsprongen hebben
gevuurd en zijn niet opnieuw vergund (de vergunningsfactoren zijn door
CDK vernietigd of geremd), dus kan zij niet opnieuw repliceren; zij
wacht tot de partner G2 bereikt en beide gaan samen de mitose in.
\end{solution}

\begin{solution}{exo:b3:cell-cycle-apoptosis:8}
De ongehechte kinetochoor is een katalysator: zij zet Mad2 doorlopend om
in zijn werkzame conformatie en maakt zo een diffundeerbare remmer van
Cdc20 die zich door de cel verspreidt en het APC/C overal uit houdt ---
één kinetochoor maakt er genoeg. Zonder Mad2 wordt het APC/C
geactiveerd zodra Cdk1 het heeft aangezet, of de chromosomen nu gehecht
zijn of niet: de anafase begint met ongehechte chromosomen, de dochters
zijn aneuploïd, en het organisme (of de cellijn) sterft aan
chromosoomverlies.
\end{solution}

\begin{solution}{exo:b3:cell-cycle-apoptosis:9}
Zonder \emph{ced-9}: cellen die zouden moeten leven sterven ---
embryonale sterfte door te veel dood. Zonder \emph{ced-3}: geen van de
$131$ sterfgevallen treedt op, de cellen overleven en differentiëren, en
de worm is levensvatbaar. Zonder beide: geen dood --- het fenotype van
\emph{ced-3}. Omdat het weghalen van de beul het effect van het
weghalen van de bewaker opheft, werkt CED-9 stroomopwaarts, door
CED-4/CED-3 inactief te houden; is hij weg, dan zijn zij werkzaam,
tenzij zij er ook niet zijn.
\end{solution}

\begin{solution}{exo:b3:cell-cycle-apoptosis:10}
Met $A = aC$ geldt $\mathrm{d}C/\mathrm{d}t = k_{s} - k_{d}aC^{2}$, een
vergelijking in één variabele met evenwichtstoestand
$C^{*} = \sqrt{k_{s}/(k_{d}a)}$ en daar een helling
$-2k_{d}aC^{*} < 0$: elke afwijking dooft eentonig uit, en één enkele
eerste-ordevariabele kan niet oscilleren. De S-vormige curve geeft twee
stabiele takken gescheiden door een verboden zone, zodat de toestand
moet springen en zich niet kan vestigen; een uitdrukkelijke vertraging
in de terugkoppeling zou op een andere manier oscillaties geven --- de
rem die te laat komt, doorschiet, enzovoort --- zoals bij veel hormonale
ritmes.
\end{solution}

\begin{solution}{exo:b3:cell-cycle-apoptosis:11}
Elke aankomende sensor wordt door een bewaker gevangen, zodat er niets
gebeurt tot alle $N$ bewakers bezet zijn, bij $t \approx N/r$; de
volgende sensoren zijn vrij, activeren Bax/Bak, waarvan de porievorming
steil met hun aantal stijgt, en de mitochondriën gaan binnen minuten
open. Een sterkere prikkel verkort de tijd alleen via $r$; boven de
drempel is de uitkomst dezelfde. Venetoclax bezet bewakers en verlaagt
zo de effectieve $N$: de tijd krimpt, en een voorgeladen cel --- al
dicht bij $N$ --- sterft meteen.
\end{solution}

\begin{solution}{exo:b3:cell-cycle-apoptosis:12}
Een overmaat Bcl-2 blokkeert de \omterm{def:b3:cell-cycle-apoptosis:pathways}{intrinsieke route}: B-cellen die zouden
moeten sterven wanneer hun groeisignalen ophouden blijven leven, en de
populatie groeit doordat de dood uitblijft, met het trage tempo waarin
zulke cellen worden gemaakt --- sluimerend, tot een tweede beschadiging
er deling bij doet. Het verlies van Rb haalt het \omterm{def:b3:cell-cycle-apoptosis:checkpoints}{restrictiepunt} weg: de
voorlopers van het netvlies doorlopen de cyclus zonder groeifactoren en
delen zich zo snel als de motor toelaat --- agressief. De twee tumoren
zijn de twee programma's die elk op zichzelf falen: een besturing van
de dood en een van de deling, waarvan het verlies van elk voor een
tumor volstaat, met de tweede sneller.
\end{solution}

\begin{solution}{pb:b3:cell-cycle-apoptosis:1}
\textbf{1.} $10^{7}\times 3500/5 = 7\times 10^{9}$ cellen per dag,
ongeveer \num{80000} per seconde.
\textbf{2.} $3500/5 = 700$ per crypte per dag; $22\times 2^{5} = 704$.
\textbf{3.} $80\times 365 \approx \num{29000}$ delingen; $\times 0.6
\approx \num{17500}$ mutaties.
\textbf{4.} Transitcellen worden binnen dagen afgestoten en nemen hun
mutaties mee; de stamcel blijft een leven lang, zodat elk van haar
mutaties behouden blijft en elke deling er meer bij doet --- traag
delen houdt het aantal klein.
\textbf{5.} De villus loopt leeg in ongeveer de $5$ dagen van zijn
normale doorvoer; de heropbouw kost een dag herstel, vijf delingen van
\qty{12}{h} en de tocht omhoog langs de villus --- vier tot vijf dagen,
waarin de barrière kaal is: de snelle transitdelingen van de darm maken
haar een vroeg slachtoffer van straling.
\textbf{6.} Een gesloten apoptotisch lijk laat niets vrij in een holte
vol bacteriën en wekt geen ontsteking op; \omterm{def:b3:cell-cycle-apoptosis:apoptosis}{necrose} zou bij elke
afgestoten cel enzymen en signalen morsen en het slijmvlies ontsteken.
\textbf{7.} $(40 - 8)/2 = \qty{16}{h}$.
\textbf{8.} $\log_{2}(40/8)\times \qty{10}{min} = 2.3\times 10 =
\qty{23}{min}$.
\textbf{9.} Periode ongeveer \qty{16.4}{h}; Cdk1 werkzaam gedurende
\qty{2.3}{\%} van de tijd.
\textbf{10.} Niet $k_{s}$ alleen, maar de tijd die bij de drempels wordt
gewacht: de somatische cel wacht bij het \omterm{def:b3:cell-cycle-apoptosis:checkpoints}{restrictiepunt} op
groeifactoren en bij G2/M op haar DNA, en de stamcel langer dan de
transitcel. Het kikkerei heeft geen \omterm{def:b3:cell-cycle-apoptosis:checkpoints}{controlepunten} en een cytoplasma dat
van alles is voorzien, zodat alleen de kale oscillator zijn periode
bepaalt.
\textbf{11.} Op de lage tak bij een cyclineniveau boven de gewone
$C_{2}$: het \omterm{def:b3:cell-cycle-apoptosis:checkpoints}{controlepunt} heeft de S-curve naar rechts geschoven door
Cdc25 te remmen, zodat de sprong uitblijft. Na het herstel komt Cdc25
vrij, zakt de drempel onder het opgehoopte \omterm{def:b3:cell-cycle-apoptosis:cdk}{cycline}, en gaat de cel
meteen de mitose in --- en daarom gaan cellen die uit een \omterm{def:b3:cell-cycle-apoptosis:checkpoints}{controlepunt}
worden vrijgelaten gelijktijdig de mitose in.
\textbf{12.} Noch \omterm{def:b3:cell-cycle-apoptosis:cdk}{cycline} B noch securine wordt afgebroken: de cel
blijft onbeperkt in de metafase met cohesine intact en Cdk1 werkzaam;
zij sterft daar of glipt er uiteindelijk uit met een niet verdeeld
genoom.
\textbf{13.} $10^{11}$ delingen per dag.
\textbf{14.} Met het \omterm{def:b3:cell-cycle-apoptosis:checkpoints}{controlepunt} $10^{11}/5\times 10^{4} = 2\times 10^{6}$
aneuploïde dochters per dag; zonder, $2\times 10^{9}$.
\textbf{15.} $2\times 10^{4}$ per dag; over $80$ jaar $6\times 10^{8}$.
\textbf{16.} $10^{9}/50 = 2\times 10^{7}$ verkeerde verdelingen per dag:
elke dag beproeft de tumor twintig miljoen nieuwe karyotypen, en de
selectie houdt die welke sneller groeien of een geneesmiddel weerstaan.
\textbf{17.} Als elke deling verkeerd verdeelt, houdt geen enkele
cellijn van het embryo een euploïd genoom en sterft het embryo.
Gedeeltelijk verlies geeft af en toe aneuploïdie in cellen die het
overleven (vooral zonder p53), een chromosomale instabiliteit die een
tumor variatie levert zonder het organisme te doden.
\textbf{18.} Een meiotische fout zet het extra chromosoom in de gameet,
zodat de zygote en elke cel die eruit voortkomt trisoom zijn; een
mitotische fout treedt in één somatische cel op en wordt alleen door
haar kloon geërfd.
\textbf{19.} $10^{11}\times \qty{1}{ng} = \qty{100}{g}$ per dag,
\qty{36}{kg} per jaar --- elk jaar ongeveer de helft van de
lichaamsmassa.
\textbf{20.} \qty{20}{g} eiwit per dag, ongeveer een derde van de
voedselinname en enkele procenten van de totale eiwitomzet van het
lichaam.
\textbf{21.} $10^{11}$ lijken bij $24$ per macrofaag per dag: $4\times
10^{9}$ macrofagen voltijds, \qty{4}{\%} van de tijd van de $10^{11}$
macrofagen.
\textbf{22.} Een uiteengevallen lijk laat proteasen, DNA, ATP en andere
alarmsignalen los die ontsteking veroorzaken en auto-immuniteit tegen
kernantigenen kunnen uitlokken; het intacte lijk meldt zich met
\omterm{def:b3:cell-cycle-apoptosis:apoptosis}{fosfatidylserine} op zijn buitenste laag (en doordat het zijn ``eet mij
niet''-signalen verliest), en lokt fagocyten met vrijgelaten
nucleotiden.
\textbf{23.} De \omterm{def:b3:cell-cycle-apoptosis:pathways}{intrinsieke route} is bij het mitochondrion geblokkeerd.
De extrinsieke weg werkt nog waar caspase-8 caspase-3 rechtstreeks
activeert, al is haar versterking via Bid verloren. Cellen hopen zich op
doordat zij niet sterven en niet doordat zij sneller delen --- trage
groei --- tot een latere beschadiging (vaak \emph{MYC}) er deling bij
doet.
\textbf{24.} De \omterm{def:b3:cell-cycle-apoptosis:apoptosis}{apoptose} kost uren, vergt ATP en doorloopt stappen die
kunnen worden geblokkeerd --- caspaseremmers, de mitochondriale drempel
--- zodat neuronen in de schemerzone te redden zijn als zij op tijd
worden behandeld; de neuronen in de kern sterven binnen minuten door
energiegebrek, zonder programma om te onderbreken.
\textbf{25.} Periode van de kale oscillator ongeveer \qty{16}{h}; zo'n $2\times
10^{4}$ aneuploïde delende cellen per dag in een normaal lichaam;
\qty{100}{g} cellen per dag door \omterm{def:b3:cell-cycle-apoptosis:apoptosis}{apoptose} hergebruikt.
\end{solution}
